\documentclass[12pt]{article}
\usepackage{amsmath, amssymb}
\usepackage[margin=1in]{geometry}
\setlength{\parskip}{6pt}
\title{QUBO Formulation: Traffic Route Assignment Problem}
\date{}
\begin{document}
\maketitle

\subsection*{Traffic Route Assignment Problem}

\paragraph{Problem Statement.}
Given a set of $N$ cars and a set of candidate routes for each car, the goal is to assign exactly one route to each car. Each route consists of a subset of road segments from a global set $S$. When multiple assigned routes share a road segment, a congestion cost is incurred. The objective is to minimize the total congestion cost, defined as the sum over all road segments of the number of pairs of selected routes that traverse that segment.

\paragraph{Decision Variables.}
Let $x_{c,r} \in \{0,1\}$ be a binary variable for each car $c \in \{0, \dots, N-1\}$ and each candidate route $r \in \{0, \dots, R_c-1\}$, where $x_{c,r} = 1$ if car $c$ is assigned route $r$, and $0$ otherwise.

\paragraph{QUBO Derivation.}
\textbf{Step 1.} The original objective is to minimize the total overlap cost:
\begin{align}
  \min \quad \sum_{s \in S} \sum_{(c,r) \in I_s} \sum_{(c',r') \in I_s, (c,r) < (c',r')} x_{c,r} x_{c',r'},
\end{align}
where $I_s$ denotes the set of car-route pairs that use road segment $s$, and the ordering $(c,r) < (c',r')$ prevents double-counting.

\textbf{Step 2.} The assignment constraints require that each car selects exactly one route:
\begin{align}
  \sum_{r=0}^{R_c-1} x_{c,r} = 1, \quad \forall c \in \{0, \dots, N-1\}.
\end{align}

\textbf{Step 3.} Each constraint is converted to a quadratic penalty term with weight $\lambda > 0$. For a fixed car $c$, the penalty is:
\begin{align}
  P_c(x) &= \lambda \left( \sum_{r=0}^{R_c-1} x_{c,r} - 1 \right)^2.
\end{align}
Expanding the square yields:
\begin{align}
  P_c(x) &= \lambda \left( \sum_{r=0}^{R_c-1} x_{c,r}^2 + 2 \sum_{0 \le r < r' \le R_c-1} x_{c,r} x_{c,r'} - 2 \sum_{r=0}^{R_c-1} x_{c,r} + 1 \right).
\end{align}
Applying the idempotent property $x_{c,r}^2 = x_{c,r}$ for binary variables simplifies the expression to:
\begin{align}
  P_c(x) &= \lambda \left( 1 - \sum_{r=0}^{R_c-1} x_{c,r} + 2 \sum_{0 \le r < r' \le R_c-1} x_{c,r} x_{c,r'} \right).
\end{align}
This is the general penalty contribution for car $c$.

\textbf{Step 4.} Combining the objective and all penalty terms gives the QUBO Hamiltonian $H(x)$:
\begin{align}
  H(x) &= \sum_{s \in S} \sum_{(c,r) \in I_s} \sum_{(c',r') \in I_s, (c,r) < (c',r')} x_{c,r} x_{c',r'} + \lambda \sum_{c=0}^{N-1} \left( 1 - \sum_{r=0}^{R_c-1} x_{c,r} + 2 \sum_{0 \le r < r' \le R_c-1} x_{c,r} x_{c,r'} \right).
\end{align}

\textbf{Step 5.} Reading off the coefficients from $H(x) = \sum_{i} Q_{i,i} x_i + \sum_{i < j} Q_{i,j} x_i x_j + c_0$, the Q-matrix entries and offset are given in closed form as:
\begin{align}
  Q_{(c,r),(c,r)} &= -\lambda, \quad \forall c, r, \\
  Q_{(c,r),(c,r')} &= 2\lambda, \quad \forall c, \, r \neq r', \\
  Q_{(c,r),(c',r')} &= \begin{cases} 
    1 & \text{if } \exists s \in S \text{ such that } (c,r) \in I_s \text{ and } (c',r') \in I_s, \\
    0 & \text{otherwise},
  \end{cases} \\
  c_0 &= \lambda N.
\end{align}

\paragraph{Penalty Weight Justification.}
Let $W_{\max}$ denote the maximum possible value of the objective function over all feasible assignments. Since each pairwise term in the objective is bounded by 1 and there are at most $\binom{N}{2}|S|$ such terms, $W_{\max} \le \binom{N}{2}|S|$. To guarantee that any infeasible assignment incurs a strictly higher energy than the optimal feasible assignment, the penalty weight must satisfy $\lambda > W_{\max}$. This ensures that the energy penalty for violating a single constraint strictly dominates any potential reduction in the objective cost. In the specific instance provided, the maximum pairwise coefficient is 1 and the maximum possible objective value is 1, so $\lambda > 1$ suffices.

\paragraph{Correctness Argument.}
Minimizing $H(x)$ over $\{0,1\}^M$ recovers the optimal route assignment because any feasible solution satisfies all constraints, yielding zero penalty energy. Consequently, the energy of a feasible solution equals its original objective cost. Any infeasible solution violates at least one constraint, incurring a penalty of at least $\lambda$. Since $\lambda$ is chosen to strictly exceed $W_{\max}$, the energy of any infeasible solution is strictly greater than the energy of the optimal feasible solution. Therefore, the global minimum of $H(x)$ must correspond to a feasible assignment that minimizes the original congestion cost.

\paragraph{Illustrative Example.}
Consider an instance with $N=2$ cars and $R_0=R_1=2$ candidate routes per car. Let the route-segment mapping be defined such that segment $s_1$ is shared by routes $(0,0)$ and $(1,0)$, and segment $s_2$ is shared by routes $(0,1)$ and $(1,1)$. Thus, $I_{s_1} = \{(0,0), (1,0)\}$ and $I_{s_2} = \{(0,1), (1,1)\}$. Using the general formulas with $\lambda = 10$:

The concrete objective is:
\begin{align}
  \text{Objective}(x) = x_{0,0} x_{1,0} + x_{0,1} x_{1,1}.
\end{align}

The concrete penalty terms (one per car) expand to:
\begin{align}
  P_0(x) &= 10 \left( 1 - x_{0,0} - x_{0,1} + 2 x_{0,0} x_{0,1} \right), \\
  P_1(x) &= 10 \left( 1 - x_{1,0} - x_{1,1} + 2 x_{1,0} x_{1,1} \right).
\end{align}

Substituting these into the general Hamiltonian yields the concrete QUBO:
\begin{align}
  H(x) &= 20 - 10 x_{0,0} - 10 x_{0,1} - 10 x_{1,0} - 10 x_{1,1} + 20 x_{0,0} x_{0,1} + 20 x_{1,0} x_{1,1} + x_{0,0} x_{1,0} + x_{0,1} x_{1,1}.
\end{align}

Enumerating the feasible assignments reveals two optimal solutions. For $x_{0,0}=1, x_{1,1}=1$ (all other variables $0$), the objective cost is $0$ and the energy is $H(x) = 20 - 10 - 10 = 0$. Similarly, $x_{0,1}=1, x_{1,0}=1$ yields energy $0$. Any assignment violating the constraint (e.g., $x_{0,0}=1, x_{0,1}=1$) incurs an energy of at least $10$, confirming the optimality of the feasible assignments.

\end{document}